(1) 在核外经典允许区要求 \(E\ge U(r)=\kappa/r\)，因此外侧转向半径 \(r_t=\kappa/E>R\)，而 \(R<r<r_t\) 是粒子经典禁阻区域。

(2) WKB 穿透率首阶为 \(P\simeq e^{-G(E)}\)，其中
\[
G(E)=\frac{2}{\hbar}\int_R^{r_t}\sqrt{2m(\kappa/r-E)}\,dr.
\]
令 \(r=r_t\sin^2\theta\)，下限 \(\theta_R=\arcsin\sqrt{R/r_t}\)。则 \(\kappa/r-E=E\cot^2\theta\)，积分化为 \(2r_t\sqrt{2mE}\int_{\theta_R}^{\pi/2}\cos^2\theta\,d\theta\)。所以
\[
G(E)=\frac{2\sqrt{2m}\kappa}{\hbar\sqrt E}
\left[
\arccos\sqrt{\frac{RE}{\kappa}}-
\sqrt{\frac{RE}{\kappa}\left(1-\frac{RE}{\kappa}\right)}
\right].
\]
其整体无量纲，是势垒作用量与 \(\hbar \) 的比值。

(3) 当 \(RE/\kappa\ll1\)，利用 \(\arccos\sqrt x=\pi/2-\sqrt x+O(x^{3/2})\) 与 \(\sqrt{x(1-x)}=\sqrt x+O(x^{3/2})\)，得
\[
G(E)=\frac{\pi\sqrt{2m}\kappa}{\hbar\sqrt E}
-\frac{4\sqrt{2m\kappa R}}{\hbar}+O(\sqrt E).
\]
主项随 \(E^{-1/2}\) 增长，故较小放出能量使穿透率指数级减小。

(4) 以 \(\nu \) 表示核内粒子单位时间尝试穿透次数，衰变率近似 \(\Gamma=\nu e^{-G}\)，对应半衰期 \(t_{1/2}=\ln2/\Gamma\)。因此
\[
\ln t_{1/2}\simeq G(E)+\ln(\ln2/\nu).
\]
在近似固定核电荷、质量和作用半径时，主要能量依赖是 \(\ln t_{1/2}\propto1/\sqrt E\)，对应 Geiger–Nuttall 型关系。但真实 \(\alpha \) 粒子预形成概率、角动量势与核力可改变前因子 \(\nu \)，WKB 仅计算此简单 Coulomb 势垒的主要指数。
